%% kapa: exam=1 round=33 year=2022 period=1 form=A subject=civil_law subject_ko=민법 seq=31 number=31
%% source: Q-Net 감정평가사 기출문제 (cst00309) artlSeq=5209997 file=raw/1008/5209997/2022년도 제33회 감정평가사 1차시험 1교시(민법,경제학,부동산학).pdf
%% conversion: kordoc 4.24.51 -> tools/gen1.py (auto); review=proofread-sonnet; flags=none
\begin{question}{31}
지역권에 관한 설명으로 옳지 않은 것은? (다툼이 있으면 판례에 따름)
\begin{choices}
  \item 통행지역권의 점유취득시효는 승역지 위에 도로를 설치하여 늘 사용하는 객관적 상태를 전제로 한다.
  \item 요역지의 공유자 중 1인이 지역권을 취득한 때에는 다른 공유자도 이를 취득한다.
  \item 요역지의 공유자 중 1인에 의한 지역권소멸시효의 중단은 다른 공유자에게는 효력이 없다.
  \item 점유로 인한 지역권취득기간의 중단은 지역권을 행사하는 모든 공유자에 대한 사유가 아니면 그 효력이 없다.
  \item 통행지역권을 시효취득한 요역지 소유자는 특별한 사정이 없는 한 승역지에 대한 도로 설치 및 사용에 의하여 승역지 소유자가 입은 손해를 보상해야 한다.
\end{choices}
\end{question}
\answer{3}
